
\documentclass[12pt,a4paper]{article}

\usepackage{amsmath,amssymb,amsthm}
\usepackage{geometry}
\usepackage{booktabs}
\usepackage{array}
\usepackage{enumitem}
\usepackage[hidelinks]{hyperref}
\usepackage{xurl}


\geometry{margin=2.5cm}
\usepackage{xepersian}
\IfFontExistsTF{B Nazanin}{\settextfont{B Nazanin}}{\settextfont{Amiri}}
\setlatintextfont{Liberation Serif}

\title{جزوه پردازش هیستوگرام تصویر\\
	\large \lr{Histogram Equalization} و \lr{Histogram Matching}}
\author{}
\date{}

\begin{document}
	
	\maketitle
	\tableofcontents
	\newpage
	
	% =========================================================
	\section{مقدمه: پردازش هیستوگرام}
	
	هیستوگرام تصویر یکی از مهم‌ترین ابزارهای تحلیل و بهبود تصویر است.
	هیستوگرام نشان می‌دهد که هر سطح شدت روشنایی چند بار در تصویر ظاهر شده است.
	
	برای یک تصویر خاکستری، محور افقی نشان‌دهنده سطح شدت و محور عمودی
	نشان‌دهنده تعداد پیکسل‌های دارای آن شدت است.
	
	اگر تصویر دارای $L$ سطح شدت باشد، آنگاه:
	
	\[
	r \in [0,L-1]
	\]
	
	که در آن:
	
	\[
	r=0
	\]
	
	نمایانگر سیاه و
	
	\[
	r=L-1
	\]
	
	نمایانگر سفید است.
	
	هدف روش‌های پردازش هیستوگرام می‌تواند یکی از موارد زیر باشد:
	
	\begin{enumerate}
		\item افزایش کنتراست تصویر؛
		\item روشن‌تر یا تیره‌تر کردن تصویر؛
		\item نزدیک کردن هیستوگرام به یک شکل مطلوب؛
		\item استخراج اطلاعات آماری از تصویر.
	\end{enumerate}
	
	
	% =========================================================
	\section{تبدیلات شدت و متغیر تصادفی شدت}
	
	شدت‌های تصویر را می‌توان به صورت یک متغیر تصادفی در بازه
	
	\[
	[0,L-1]
	\]
	
	در نظر گرفت.
	
	فرض کنید:
	
	\[
	r
	\]
	
	شدت پیکسل در تصویر ورودی باشد و
	
	\[
	s
	\]
	
	شدت متناظر در تصویر خروجی باشد.
	
	یک تبدیل شدت به صورت زیر تعریف می‌شود:
	
	\[
	s=T(r)
	\]
	
	یا، در فرم نرمال‌شده‌ای که در برخی منابع استفاده می‌شود:
	
	\[
	s=T(r),\qquad 0\le r\le L-1.
	\]
	
	برای تبدیل شدت معمولاً دو شرط در نظر گرفته می‌شود.
	
	\subsection{شرط اول: یکنواختی}
	
	تابع تبدیل باید یک تابع صعودی یکنواخت باشد:
	
	\[
	T(r_1)\le T(r_2)
	\qquad
	\text{اگر}
	\qquad
	r_1<r_2.
	\]
	
	این شرط مانع وارونه شدن ترتیب شدت‌ها می‌شود.
	
	یعنی اگر یک پیکسل از پیکسل دیگر روشن‌تر باشد، پس از تبدیل نیز
	نباید تیره‌تر از آن شود.
	
	\subsection{شرط دوم: حفظ محدوده شدت}
	
	باید داشته باشیم:
	
	\[
	0\le T(r)\le L-1.
	\]
	
	بنابراین خروجی همچنان در همان محدوده شدت تصویر قرار می‌گیرد.
	
	
	% =========================================================
	\section{یک‌به‌یک و چندبه‌یک بودن تبدیل}
	
	اگر $T(r)$ فقط یکنواخت افزایشی باشد، ممکن است چند مقدار ورودی مختلف
	به یک مقدار خروجی نگاشت شوند.
	
	برای مثال:
	
	\[
	r_1<r_2<r_3
	\]
	
	ولی:
	
	\[
	T(r_1)=T(r_2)=T(r_3).
	\]
	
	در این حالت تبدیل از نوع
	
	\[
	\text{\lr{many-to-one}}
	\]
	
	است.
	
	این موضوع برای تبدیل $r$ به $s$ مشکلی ایجاد نمی‌کند، اما اگر بخواهیم
	از $s$ به $r$ برگردیم، ممکن است نتوانیم مقدار اصلی $r$ را به صورت
	منحصربه‌فرد مشخص کنیم.
	
	برای جلوگیری از این مشکل، تابع باید \textbf{\lr{strictly monotonic increasing}}
	باشد:
	
	\[
	r_1<r_2
	\quad\Longrightarrow\quad
	T(r_1)<T(r_2).
	\]
	
	در این صورت تابع یک‌به‌یک است و معکوس آن نیز یکتا خواهد بود:
	
	\[
	r=T^{-1}(s).
	\]
	
	
	% =========================================================
	\section{تبدیل متغیر تصادفی شدت}
	
	فرض کنید $r$ دارای تابع چگالی احتمال
	
	\[
	p_r(r)
	\]
	
	باشد.
	
	پس از اعمال تبدیل
	
	\[
	s=T(r)
	\]
	
	متغیر جدید $s$ دارای تابع چگالی احتمال
	
	\[
	p_s(s)
	\]
	
	خواهد بود.
	
	رابطه بنیادی تبدیل متغیرهای تصادفی چنین است:
	
	\[
	\boxed{
		p_s(s)
		=
		p_r(r)
		\left|
		\frac{dr}{ds}
		\right|
	}
	\tag{1}
	\]
	
	یا به صورت معادل:
	
	\[
	\boxed{
		p_s(s)
		=
		p_r(r)
		\left|
		\frac{1}{ds/dr}
		\right|
	}
	\tag{2}
	\]
	
	این رابطه اساس ریاضی \lr{Histogram Equalization} است.
	
	
	% =========================================================
	\section{\lr{Histogram Equalization} در حالت پیوسته}
	
	تبدیل بسیار مهم زیر را تعریف می‌کنیم:
	
	\[
	\boxed{
		s=T(r)
		=
		(L-1)
		\int_0^r p_r(w)\,dw
	}
	\tag{3}
	\]
	
	در این رابطه، $w$ یک متغیر کمکی برای انتگرال است.
	
	انتگرال
	
	\[
	\int_0^r p_r(w)\,dw
	\]
	
	در واقع تابع توزیع تجمعی یا
	
	\[
	\text{\lr{CDF}}
	\]
	
	متغیر تصادفی $r$ است.
	
	بنابراین می‌توان نوشت:
	
	\[
	F_r(r)=\int_0^r p_r(w)\,dw
	\]
	
	و در نتیجه:
	
	\[
	\boxed{
		s=(L-1)F_r(r)
	}
	\tag{4}
	\]
	
	این مهم‌ترین رابطه \lr{Histogram Equalization} است.
	
	
	% =========================================================
	\section{چرا تبدیل فوق یکنواخت است؟}
	
	از آنجا که \lr{PDF} همواره نامنفی است:
	
	\[
	p_r(r)\ge 0,
	\]
	
	با افزایش $r$، مساحت زیر منحنی \lr{PDF} نمی‌تواند کاهش پیدا کند.
	
	بنابراین:
	
	\[
	\int_0^r p_r(w)\,dw
	\]
	
	یک تابع صعودی یکنواخت است.
	
	پس:
	
	\[
	T(r)
	=
	(L-1)
	\int_0^r p_r(w)\,dw
	\]
	
	نیز صعودی یکنواخت است.
	
	همچنین چون:
	
	\[
	\int_0^{L-1}p_r(w)\,dw=1,
	\]
	
	داریم:
	
	\[
	T(L-1)
	=
	(L-1).
	\]
	
	و:
	
	\[
	T(0)=0.
	\]
	
	بنابراین:
	
	\[
	0\le T(r)\le L-1.
	\]
	
	پس شرایط لازم برای تبدیل شدت برقرار است.
	
	
	% =========================================================
	\section{اثبات یکنواخت شدن \lr{PDF}}
	
	از رابطه
	
	\[
	s=T(r)
	\]
	
	داریم:
	
	\[
	\frac{ds}{dr}
	=
	\frac{dT(r)}{dr}.
	\]
	
	با استفاده از رابطه (3):
	
	\[
	\frac{ds}{dr}
	=
	(L-1)
	\frac{d}{dr}
	\left[
	\int_0^r p_r(w)\,dw
	\right].
	\]
	
	طبق قضیه اساسی حسابان:
	
	\[
	\frac{d}{dr}
	\left[
	\int_0^r p_r(w)\,dw
	\right]
	=
	p_r(r).
	\]
	
	پس:
	
	\[
	\boxed{
		\frac{ds}{dr}
		=
		(L-1)p_r(r)
	}
	\tag{5}
	\]
	
	بنابراین:
	
	\[
	\frac{dr}{ds}
	=
	\frac{1}{(L-1)p_r(r)}.
	\tag{6}
	\]
	
	اکنون رابطه تبدیل \lr{PDF} را داریم:
	
	\[
	p_s(s)
	=
	p_r(r)
	\left|
	\frac{dr}{ds}
	\right|.
	\]
	
	با جایگذاری رابطه (6):
	
	\[
	p_s(s)
	=
	p_r(r)
	\frac{1}{(L-1)p_r(r)}.
	\]
	
	بنابراین:
	
	\[
	\boxed{
		p_s(s)=\frac{1}{L-1}
	}
	\tag{7}
	\]
	
	برای:
	
	\[
	0\le s\le L-1.
	\]
	
	این دقیقاً یک \lr{PDF} یکنواخت است.
	
	بنابراین:
	
	\[
	\boxed{
		\text{\lr{Histogram Equalization}}
		\Longrightarrow
		\text{\lr{Uniform PDF}}
	}
	\]
	
	در حالت پیوسته، مستقل از شکل \lr{PDF} اولیه، خروجی دارای \lr{PDF} یکنواخت
	خواهد بود.
	
	
	% =========================================================
	\section{مثال پیوسته}
	
	فرض کنید \lr{PDF} شدت تصویر به صورت زیر باشد:
	
	\[
	p_r(r)=
	\begin{cases}
		\dfrac{2r}{(L-1)^2},
		&
		0\le r\le L-1,
		\\[6pt]
		0,
		&
		\text{\lr{otherwise}}.
	\end{cases}
	\tag{8}
	\]
	
	با استفاده از رابطه \lr{Histogram Equalization} داریم:
	
	\[
	s
	=
	(L-1)
	\int_0^r
	\frac{2w}{(L-1)^2}
	\,dw.
	\]
	
	انتگرال را محاسبه می‌کنیم:
	
	\[
	s
	=
	(L-1)
	\frac{1}{(L-1)^2}
	\left[w^2\right]_0^r.
	\]
	
	پس:
	
	\[
	s
	=
	\frac{r^2}{L-1}.
	\]
	
	بنابراین:
	
	\[
	\boxed{
		s=\frac{r^2}{L-1}
	}
	\tag{9}
	\]
	
	یعنی برای \lr{equalization} این تصویر، شدت $r$ را به صورت مربع آن
	تبدیل و سپس بر $L-1$ تقسیم می‌کنیم.
	
	از طرف دیگر:
	
	\[
	r=\sqrt{s(L-1)}.
	\]
	
	همچنین:
	
	\[
	\frac{ds}{dr}
	=
	\frac{2r}{L-1}.
	\]
	
	بنابراین:
	
	\[
	\frac{dr}{ds}
	=
	\frac{L-1}{2r}.
	\]
	
	اکنون:
	
	\[
	p_s(s)
	=
	p_r(r)
	\left|\frac{dr}{ds}\right|.
	\]
	
	پس:
	
	\[
	p_s(s)
	=
	\frac{2r}{(L-1)^2}
	\frac{L-1}{2r}.
	\]
	
	در نتیجه:
	
	\[
	\boxed{
		p_s(s)=\frac{1}{L-1}
	}
	\]
	
	که همان \lr{PDF} یکنواخت است.
	
	
	% =========================================================
	\section{\lr{Histogram Equalization} در تصاویر دیجیتال}
	
	در تصویر دیجیتال، شدت‌ها گسسته هستند.
	
	فرض کنید:
	
	\[
	r_k,\qquad k=0,1,\ldots,L-1
	\]
	
	سطوح شدت تصویر باشند.
	
	اگر تصویر دارای ابعاد
	
	\[
	M\times N
	\]
	
	باشد، تعداد کل پیکسل‌ها برابر است با:
	
	\[
	MN.
	\]
	
	فرض کنید:
	
	\[
	n_k
	\]
	
	تعداد پیکسل‌هایی باشد که شدت آنها برابر $r_k$ است.
	
	بنابراین احتمال تقریبی مشاهده شدت $r_k$ برابر است با:
	
	\[
	\boxed{
		p_r(r_k)
		=
		\frac{n_k}{MN}
	}
	\tag{10}
	\]
	
	این همان \lr{Histogram} نرمال‌شده تصویر است.
	
	همچنین:
	
	\[
	\sum_{k=0}^{L-1}n_k=MN.
	\]
	
	بنابراین:
	
	\[
	\sum_{k=0}^{L-1}p_r(r_k)=1.
	\]
	
	
	% =========================================================
	\section{رابطه گسسته \lr{Histogram Equalization}}
	
	نسخه گسسته رابطه (3) به صورت زیر است:
	
	\[
	\boxed{
		s_k
		=
		T(r_k)
		=
		(L-1)
		\sum_{j=0}^{k}
		p_r(r_j)
	}
	\tag{11}
	\]
	
	با جایگذاری رابطه (10):
	
	\[
	\boxed{
		s_k
		=
		(L-1)
		\sum_{j=0}^{k}
		\frac{n_j}{MN}
	}
	\tag{12}
	\]
	
	این رابطه اصلی \lr{Histogram Equalization} برای تصویر دیجیتال است.
	
	یعنی برای هر سطح شدت $r_k$:
	
	\begin{enumerate}
		\item تعداد پیکسل‌های آن شدت را پیدا می‌کنیم؛
		\item احتمال آن را محاسبه می‌کنیم؛
		\item احتمال‌ها را از شدت صفر تا $r_k$ جمع می‌کنیم؛
		\item حاصل را در $L-1$ ضرب می‌کنیم؛
		\item نتیجه را به نزدیک‌ترین عدد صحیح گرد می‌کنیم.
	\end{enumerate}
	
	
	% =========================================================
	\section{مثال: تصویر 3 بیتی}
	
	فرض کنید:
	
	\[
	L=8
	\]
	
	باشد.
	
	پس شدت‌ها عبارت‌اند از:
	
	\[
	0,1,2,3,4,5,6,7.
	\]
	
	فرض کنید:
	
	\[
	M=N=64.
	\]
	
	پس:
	
	\[
	MN=4096.
	\]
	
	جدول هیستوگرام ورودی:
	
	\[
	\begin{array}{c|c|c}
		r_k & n_k & p_r(r_k)=n_k/4096\\
		\hline
		0&790&0.19\\
		1&1023&0.25\\
		2&850&0.21\\
		3&656&0.16\\
		4&329&0.08\\
		5&245&0.06\\
		6&122&0.03\\
		7&81&0.02
	\end{array}
	\]
	
	از رابطه:
	
	\[
	s_k
	=
	7
	\sum_{j=0}^{k}p_r(r_j)
	\]
	
	استفاده می‌کنیم.
	
	برای $r_0=0$:
	
	\[
	s_0=7(0.19)=1.33.
	\]
	
	برای $r_1=1$:
	
	\[
	s_1
	=
	7(0.19+0.25)
	=
	3.08.
	\]
	
	برای $r_2=2$:
	
	\[
	s_2
	=
	7(0.19+0.25+0.21)
	=
	4.55.
	\]
	
	به همین ترتیب:
	
	\[
	\begin{aligned}
		s_0&=1.33,\\
		s_1&=3.08,\\
		s_2&=4.55,\\
		s_3&=5.67,\\
		s_4&=6.23,\\
		s_5&=6.65,\\
		s_6&=6.86,\\
		s_7&=7.00.
	\end{aligned}
	\]
	
	با گرد کردن:
	
	\[
	\boxed{
		\begin{aligned}
			0&\rightarrow1\\
			1&\rightarrow3\\
			2&\rightarrow5\\
			3&\rightarrow6\\
			4&\rightarrow6\\
			5&\rightarrow7\\
			6&\rightarrow7\\
			7&\rightarrow7
		\end{aligned}
	}
	\tag{13}
	\]
	
	بنابراین فقط پنج سطح خروجی مختلف ایجاد شده‌اند:
	
	\[
	1,3,5,6,7.
	\]
	
	
	% =========================================================
	\section{نکته مهم درباره \lr{Histogram} گسسته}
	
	در حالت پیوسته ثابت کردیم:
	
	\[
	p_s(s)=\frac{1}{L-1}.
	\]
	
	اما در تصویر دیجیتال، این نتیجه به صورت دقیق تضمین نمی‌شود.
	
	علت:
	
	\begin{enumerate}
		\item شدت‌ها گسسته‌اند؛
		\item تعداد پیکسل‌ها محدود است؛
		\item مقدارهای $s_k$ باید گرد شوند؛
		\item سطح شدت جدیدی در فرآیند ایجاد نمی‌شود.
	\end{enumerate}
	
	بنابراین \lr{Histogram Equalization} گسسته الزاماً یک هیستوگرام کاملاً
	صاف تولید نمی‌کند.
	
	هدف اصلی، \textbf{گسترش مناسب توزیع شدت‌ها} و افزایش کنتراست است.
	
	
	% =========================================================
	\section{چرا \lr{Histogram Equalization} کنتراست را افزایش می‌دهد؟}
	
	فرض کنید تعداد زیادی پیکسل در قسمت تاریک تصویر متمرکز باشند.
	
	در این حالت:
	
	\[
	p_r(r)
	\]
	
	در مقادیر پایین زیاد است.
	
	\lr{CDF} در این ناحیه سریع افزایش پیدا می‌کند:
	
	\[
	F_r(r)=\int_0^r p_r(w)\,dw.
	\]
	
	بنابراین:
	
	\[
	s=(L-1)F_r(r)
	\]
	
	نیز سریع افزایش می‌یابد.
	
	در نتیجه مقادیر شدت تاریک به دامنه بزرگ‌تری از شدت‌ها نگاشت می‌شوند.
	
	پس:
	
	\[
	\boxed{
		\text{مقادیر متراکم}
		\rightarrow
		\text{پراکنده‌تر شدن}
		\rightarrow
		\text{افزایش کنتراست}
	}
	\]
	
	
	% =========================================================
	

% =========================================================
\section{محدودیت‌های \lr{Histogram Equalization} و مسیرهای توسعه معتبر}

متعادل‌سازی سراسری هیستوگرام به دلیل سادگی، هزینه محاسباتی کم و تفسیر روشن مبتنی بر
تابع توزیع تجمعی، نقطه شروع بسیار خوبی برای بهبود کنتراست است. با این حال، یک تبدیل
سراسری الزاماً برای همه تصاویر بهترین پاسخ نیست. در عمل سه مسئله مهم دیده می‌شود:
تغییر میانگین روشنایی، تقویت بیش از حد کنتراست یا نویز، و ناتوانی یک هیستوگرام سراسری
در سازگاری با نواحی دارای روشنایی‌های متفاوت. این محدودیت‌ها در ادبیات کلاسیک و جدید
به شکل‌های مختلف بررسی شده‌اند \cite{kim1997,pizer1987,stark2000,arici2009}.

\subsection{حفظ روشنایی با \lr{Brightness-Preserving Bi-Histogram Equalization}}

در روش \lr{BBHE} که توسط \lr{Kim} پیشنهاد شد، هیستوگرام ورودی در حوالی میانگین شدت
به دو زیرهیستوگرام تقسیم می‌شود و هر بخش در بازه مخصوص خود متعادل می‌گردد
\cite{kim1997}. اگر میانگین شدت ورودی برابر

\[
\mu_X=\sum_{k=0}^{L-1} r_k p_r(r_k)
\]

باشد، یک آستانه گسسته را می‌توان به صورت

\[
 m=\left\lfloor \mu_X \right\rfloor
\]

در نظر گرفت. دو مجموعه شدت عبارت‌اند از

\[
\mathcal{X}_{L}=\{0,\ldots,m\},
\qquad
\mathcal{X}_{U}=\{m+1,\ldots,L-1\}.
\]

برای هر بخش، ابتدا احتمال‌ها در همان زیرمجموعه دوباره نرمال می‌شوند. برای نمونه، در
بخش پایین داریم

\[
 p_L(r_k)=
 \frac{p_r(r_k)}{\sum_{j=0}^{m}p_r(r_j)},
 \qquad 0\le k\le m,
\]

و تابع تجمعی شرطی آن برابر است با

\[
 F_L(r_k)=\sum_{j=0}^{k}p_L(r_j).
\]

یک نگاشت آموزشی سازگار با ایده \lr{BBHE} برای نیمه پایین و بالا چنین نوشته می‌شود:

\[
 s_k=
 \begin{cases}
 m\,F_L(r_k), & 0\le r_k\le m,\\[4pt]
 (m+1)+(L-m-2)F_U(r_k), & m<r_k\le L-1.
 \end{cases}
\]

نکته کلیدی آن است که دو نیمه به کل بازه شدت کشیده نمی‌شوند؛ در نتیجه تغییر روشنایی
میانگین معمولاً از متعادل‌سازی سراسری کمتر است. این روش تضمین نمی‌کند که میانگین خروجی
دقیقاً با ورودی برابر باشد، اما در مقاله اصلی نشان داده شده است که حفظ روشنایی به شکل
معناداری بهبود می‌یابد \cite{kim1997}.

\subsection{متعادل‌سازی تطبیقی و کنترل تقویت نویز}

در \lr{Adaptive Histogram Equalization (AHE)} به جای یک هیستوگرام برای کل تصویر، برای
همسایگی هر نقطه یک هیستوگرام محلی ساخته می‌شود. اگر پنجره‌ای با مرکز $(x,y)$ داشته
باشیم، هیستوگرام محلی را می‌توان به صورت

\[
 h_{x,y}(k)=
 \#\{(u,v)\in\Omega_{x,y}: I(u,v)=k\}
\]

تعریف کرد. سپس همان ایده تابع تجمعی، این بار بر اساس $h_{x,y}$، یک نگاشت محلی می‌سازد.
مزیت این روش آن است که ناحیه تاریک و روشن تصویر می‌توانند نگاشت‌های متفاوتی دریافت
کنند. با این حال، \lr{Pizer et al.} نشان دادند که \lr{AHE} می‌تواند در نواحی تقریباً
همگن نویز را بیش از حد تقویت کند و نسخه بریده‌شده را برای کنترل این مسئله مطرح کردند
\cite{pizer1987}.

در یک صورت ساده، اگر حد برش برابر $\tau$ باشد، هر ستون هیستوگرام به

\[
 \widetilde h(k)=\min\{h(k),\tau\}
\]

محدود می‌شود. جرم حذف‌شده

\[
 E=\sum_{k=0}^{L-1}\bigl(h(k)-\widetilde h(k)\bigr)
\]

می‌تواند بین خانه‌های هیستوگرام بازتوزیع شود. این همان ایده محوری است که در پیاده‌سازی
رایج \lr{CLAHE} برای جلوگیری از رشد شدید شیب تابع تجمعی محلی استفاده می‌شود. بسته کد
همراه این جزوه یک نمونه اجرایی از \lr{CLAHE} دارد.

\subsection{از متعادل‌سازی کامل تا تقویت کنترل‌شده}

\lr{Stark} نشان داد که با تعمیم تابع انباشت محلی می‌توان خانواده‌ای از تبدیلات ساخت که
بین «عدم تغییر تصویر» و «متعادل‌سازی تطبیقی کامل» حرکت کنند \cite{stark2000}. این دیدگاه
از نظر طراحی مهم است: هدف لزوماً یکنواخت کردن هرچه بیشتر هیستوگرام نیست، بلکه یافتن
شدت مناسبِ تقویت برای کاربرد موردنظر است.

همین ایده در چارچوب‌های جدیدتر به صورت یک مسئله بهینه‌سازی دیده می‌شود. برای مثال،
\lr{Arici et al.} یک چارچوب اصلاح هیستوگرام ارائه کردند که در آن می‌توان میان افزایش
کنتراست، طبیعی ماندن ظاهر، مقاومت در برابر نویز و حفظ روشنایی سازش ایجاد کرد
\cite{arici2009}. برای فهم این فلسفه، یک قالب آموزشی مناسب چنین است:

\[
 q^{\star}
 =
 \arg\min_{q\in\Delta_L}
 \left[
 \alpha D(q,u)
 +(1-\alpha)D(q,p)
 +\lambda R(q)
 \right],
\]

که در آن $p$ هیستوگرام نرمال‌شده ورودی، $u$ توزیع یکنواخت، $q$ هیستوگرام هدف،
$D(\cdot,\cdot)$ یک فاصله یا واگرایی مناسب و $R(q)$ یک جمله منظم‌ساز است. این رابطه
عین تابع هزینه مقاله یادشده نیست؛ بلکه یک بازنویسی آموزشی برای نشان دادن مفهوم سازش
میان «نزدیکی به یکنواخت» و «نزدیکی به تصویر اصلی» است. هرچه $\alpha$ بزرگ‌تر باشد،
تمایل به افزایش کنتراست بیشتر می‌شود و هرچه وزن بخش دوم بیشتر باشد، تغییر شکل
هیستوگرام محدودتر می‌شود.

\subsection{نکته عملی برای تصاویر رنگی}

اگر سه کانال رنگی به صورت مستقل متعادل شوند، نسبت بین کانال‌ها می‌تواند تغییر کند و
رنگ‌های غیرطبیعی ایجاد شود. یک راه عملی این است که تصویر به فضایی شامل مؤلفه روشنایی
تبدیل شود، عملیات هیستوگرام فقط روی روشنایی انجام گیرد و سپس تصویر به فضای اولیه
برگردانده شود. در کد همراه، برای تصویر رنگی از مؤلفه روشنایی فضای \lr{Lab} استفاده شده
است تا تغییر رنگ کاهش یابد.

% =========================================================
\section{نگاه نظریه اطلاعات به هیستوگرام}

هیستوگرام نرمال‌شده را می‌توان یک توزیع احتمال گسسته دانست. همین مشاهده اجازه می‌دهد
ابزارهای نظریه اطلاعات مستقیماً روی تصویر اعمال شوند. تعریف کلاسیک آنتروپی از کار
\lr{Shannon} می‌آید \cite{shannon1948}. اگر

\[
 p_k=p_r(r_k),
 \qquad
 \sum_{k=0}^{L-1}p_k=1,
\]

باشد، آنتروپی شدت تصویر برابر است با

\[
 \boxed{
 H(R)=-\sum_{k=0}^{L-1}p_k\log_2 p_k
 }
\]

که جمله‌های دارای $p_k=0$ برابر صفر در نظر گرفته می‌شوند. واحد این کمیت \lr{bit} است.
اگر تنها یک سطح شدت وجود داشته باشد، $H(R)=0$ و اگر احتمال همه سطوح برابر باشد،
آنتروپی به بیشینه خود می‌رسد.

\subsection{چرا توزیع یکنواخت بیشترین آنتروپی را دارد؟}

برای توزیع یکنواخت

\[
 u_k=\frac{1}{L}
\]

واگرایی \lr{Kullback--Leibler} بین $p$ و $u$ برابر است با

\[
 D_{\mathrm{KL}}(p\Vert u)
 =
 \sum_{k=0}^{L-1}p_k
 \log_2\frac{p_k}{u_k}.
\]

بر اساس نامنفی بودن این واگرایی \cite{kullback1951} داریم

\[
 \begin{aligned}
 D_{\mathrm{KL}}(p\Vert u)
 &=\sum_k p_k\log_2 p_k+\log_2 L\\
 &=\log_2L-H(R)\ge 0.
 \end{aligned}
\]

پس

\[
 \boxed{H(R)\le \log_2 L}
\]

و برابری دقیقاً هنگامی رخ می‌دهد که $p_k=1/L$ برای همه $k$ باشد. این نتیجه یک پیوند
مستقیم میان نظریه اطلاعات و متعادل‌سازی هیستوگرام می‌سازد: در مدل ایده‌آل پیوسته،
\lr{Histogram Equalization} خروجی یکنواخت ایجاد می‌کند؛ بنابراین از دید توزیع شدت،
به حالت بیشینه‌آنتروپی نزدیک می‌شود. در حالت دیجیتال، به علت گرد کردن، تعداد محدود
پیکسل‌ها و نگاشت چندبه‌یک، هیستوگرام خروجی عموماً کاملاً یکنواخت و در نتیجه دقیقاً
بیشینه‌آنتروپی نیست.

این هویت همچنین یک معیار ساده برای فاصله از یکنواختی می‌دهد:

\[
 \boxed{
 D_{\mathrm{KL}}(p\Vert u)=\log_2L-H(R)
 }
\]

بنابراین اگر دو تصویر دارای تعداد سطح شدت یکسان باشند، افزایش آنتروپی هیستوگرام
معادل کاهش این واگرایی نسبت به توزیع یکنواخت است.

\subsection{مورد اول استخراج اطلاعات: آنتروپی چه چیزی را می‌گوید و چه چیزی را نمی‌گوید؟}

آنتروپی هیستوگرام میزان تنوع آماری شدت‌ها را خلاصه می‌کند و در کارهای مبتنی بر
هیستوگرام به عنوان یکی از معیارهای غنای توزیع به کار رفته است. برای نمونه،
\lr{Zhuang and Guan} از آنتروپی برای تقسیم تطبیقی زیرهیستوگرام‌ها استفاده کردند
\cite{zhuang2018}. با این حال باید در تفسیر آن محتاط بود: اگر مکان پیکسل‌های یک تصویر
را به طور تصادفی جابه‌جا کنیم ولی تعداد هر سطح شدت را ثابت نگه داریم، هیستوگرام و
در نتیجه $H(R)$ تغییر نمی‌کند، در حالی که ساختار فضایی تصویر ممکن است کاملاً نابود
شود.

این محدودیت دقیقاً نشان می‌دهد که «آنتروپی هیستوگرام» مترادف «جزئیات ادراکی» نیست.
\lr{Celik} برای جبران همین کمبود، آنتروپی مکانی را وارد مسئله افزایش کنتراست کرد و
به جای تکیه صرف بر فراوانی شدت‌ها، توزیع مکانی پیکسل‌های هر سطح خاکستری را نیز در
نظر گرفت \cite{celik2014}. بنابراین برای ارزیابی علمی یک روش، آنتروپی بهتر است در کنار
معیارهای ساختاری، روشنایی و مشاهده بصری استفاده شود.

\subsection{مورد دوم استخراج اطلاعات: هیستوگرام مشترک و \lr{Mutual Information}}

اگر دو تصویر $A$ و $B$ از یک صحنه، دو حسگر، یا دو زمان متفاوت داشته باشیم، می‌توان
به جای یک هیستوگرام تک‌متغیره از هیستوگرام مشترک استفاده کرد. اگر
$h_{AB}(i,j)$ تعداد پیکسل‌هایی باشد که همزمان شدت $i$ در تصویر $A$ و شدت $j$ در تصویر
$B$ دارند، احتمال مشترک به صورت

\[
 p_{ij}
 =
 \frac{h_{AB}(i,j)}{N}
\]

تعریف می‌شود. حاشیه‌ها برابرند با

\[
 p_i=\sum_j p_{ij},
 \qquad
 p_j=\sum_i p_{ij}.
\]

اطلاعات متقابل دو تصویر برابر است با

\[
 \boxed{
 I(A;B)=
 \sum_{i,j}p_{ij}
 \log_2\frac{p_{ij}}{p_i p_j}
 }
\]

و شکل معادل آن

\[
 I(A;B)=H(A)+H(B)-H(A,B)
\]

است. اگر دو متغیر مستقل باشند، $p_{ij}=p_i p_j$ و در نتیجه $I(A;B)=0$ می‌شود.
هرچه وابستگی آماری بیشتر باشد، اطلاعات متقابل نیز افزایش می‌یابد. \lr{Viola and Wells}
این ایده را به صورت مؤثر برای هم‌ترازی تصویر با بیشینه‌سازی اطلاعات متقابل به کار
بردند \cite{viola1997}. اهمیت این مثال برای بحث فعلی آن است که «اطلاعات» فقط از شکل
یک هیستوگرام استخراج نمی‌شود؛ هیستوگرام مشترک می‌تواند رابطه آماری میان دو تصویر را
نمایان کند.

% =========================================================
\section{معیارهای کمی ساده برای ارزیابی عملیات هیستوگرام}

برای مقایسه چند روش، اتکا به یک معیار کافی نیست. چند کمیت مکمل که در کد همراه نیز
محاسبه می‌شوند عبارت‌اند از:

\begin{enumerate}
 \item آنتروپی شدت $H(R)$ برای اندازه‌گیری تنوع توزیع شدت؛
 \item فاصله از توزیع یکنواخت با $D_{\mathrm{KL}}(p\Vert u)$؛
 \item اختلاف مطلق میانگین روشنایی، موسوم به \lr{AMBE}:
 \[
 \operatorname{AMBE}=|\mu_{\mathrm{in}}-\mu_{\mathrm{out}}|;
 \]
 \item انحراف معیار شدت به عنوان یک شاخص ساده از گستردگی کنتراست:
 \[
 \sigma=\sqrt{\sum_k(r_k-\mu)^2p_k}.
 \]
\end{enumerate}

کاهش \lr{AMBE} برای کاربردهایی که حفظ روشنایی مهم است مطلوب است، اما به تنهایی کیفیت
تصویر را تضمین نمی‌کند. همچنین آنتروپی بالا ممکن است همراه با نویز بالا باشد. بنابراین
خروجی مطلوب معمولاً حاصل یک سازش بین کنتراست، روشنایی، نویز و ساختار مکانی است؛ همین
موضوع انگیزه روش‌های حفظ روشنایی، تطبیقی و بهینه‌سازی‌شده است
\cite{kim1997,pizer1987,arici2009}.


\section{\lr{Histogram Matching} یا \lr{Histogram Specification}}
	
	\lr{Histogram Equalization} همیشه بهترین انتخاب نیست.
	
	گاهی می‌خواهیم \lr{Histogram} خروجی شکل خاصی داشته باشد.
	
	در این حالت از:
	
	\[
	\boxed{
		\text{\lr{Histogram Matching}}
	}
	\]
	
	یا:
	
	\[
	\boxed{
		\text{\lr{Histogram Specification}}
	}
	\]
	
	استفاده می‌کنیم.
	
	در این روش، به جای اینکه بگوییم:
	
	\[
	\text{\lr{Histogram} خروجی باید \lr{Uniform} باشد}
	\]
	
	می‌گوییم:
	
	\[
	\text{\lr{Histogram} خروجی باید شکل مشخصی داشته باشد.}
	\]
	
	
	% =========================================================
	\section{متغیر شدت خروجی مطلوب}
	
	فرض کنید:
	
	\[
	r
	\]
	
	شدت ورودی و
	
	\[
	z
	\]
	
	شدت مطلوب خروجی باشد.
	
	\lr{PDF} ورودی:
	
	\[
	p_r(r)
	\]
	
	و \lr{PDF} مطلوب:
	
	\[
	p_z(z)
	\]
	
	است.
	
	ابتدا تصویر ورودی را \lr{Equalize} می‌کنیم:
	
	\[
	\boxed{
		s=T(r)
		=
		(L-1)
		\int_0^r p_r(w)\,dw
	}
	\tag{14}
	\]
	
	بنابراین $s$ دارای \lr{PDF} یکنواخت خواهد بود.
	
	
	% =========================================================
	\section{ساخت تابع تبدیل خروجی مطلوب}
	
	اکنون برای \lr{PDF} مطلوب $p_z(z)$ تابع زیر را تعریف می‌کنیم:
	
	\[
	\boxed{
		G(z)
		=
		(L-1)
		\int_0^z p_z(v)\,dv
	}
	\tag{15}
	\]
	
	در این رابطه $v$ متغیر کمکی انتگرال است.
	
	چون $G(z)$ نیز \lr{CDF} نرمال‌شده \lr{PDF} مطلوب است، داریم:
	
	\[
	G(z)=s.
	\]
	
	از طرف دیگر:
	
	\[
	s=T(r).
	\]
	
	بنابراین:
	
	\[
	\boxed{
		G(z)=T(r)
	}
	\tag{16}
	\]
	
	در نتیجه:
	
	\[
	\boxed{
		z=G^{-1}(s)
	}
	\tag{17}
	\]
	
	و چون:
	
	\[
	s=T(r),
	\]
	
	می‌توان مستقیماً نوشت:
	
	\[
	\boxed{
		z=G^{-1}[T(r)]
	}
	\tag{18}
	\]
	
	این رابطه اساس \lr{Histogram Specification} است.
	
	
	% =========================================================
	\section{تفسیر ریاضی \lr{Histogram Matching}}
	
	کل فرآیند را می‌توان به صورت زیر دید:
	
	\[
	r
	\xrightarrow{\ T\ }
	s
	\xrightarrow{\ G^{-1}\ }
	z
	\]
	
	مرحله اول:
	
	\[
	r\rightarrow s
	\]
	
	هیستوگرام ورودی را \lr{Equalize} می‌کند.
	
	مرحله دوم:
	
	\[
	s\rightarrow z
	\]
	
	توزیع یکنواخت را به توزیع مطلوب تبدیل می‌کند.
	
	در نتیجه:
	
	\[
	\boxed{
		r
		\rightarrow
		s=T(r)
		\rightarrow
		z=G^{-1}(s)
	}
	\]
	
	
	% =========================================================
	\section{الگوریتم \lr{Histogram Specification} در حالت پیوسته}
	
	\begin{enumerate}
		\item \lr{PDF} تصویر ورودی را به دست آورید:
		
		\[
		p_r(r).
		\]
		
		\item با استفاده از آن تابع $T(r)$ را بسازید:
		
		\[
		T(r)
		=
		(L-1)\int_0^r p_r(w)\,dw.
		\]
		
		\item \lr{PDF} مطلوب را مشخص کنید:
		
		\[
		p_z(z).
		\]
		
		\item تابع $G(z)$ را بسازید:
		
		\[
		G(z)
		=
		(L-1)\int_0^z p_z(v)\,dv.
		\]
		
		\item معکوس $G$ را پیدا کنید:
		
		\[
		G^{-1}(s).
		\]
		
		\item برای هر پیکسل:
		
		\[
		z=G^{-1}[T(r)].
		\]
	\end{enumerate}
	
	
	% =========================================================
	\section{فرم گسسته \lr{Histogram Specification}}
	
	در تصویر دیجیتال، به جای \lr{PDF} از \lr{Histogram} نرمال‌شده استفاده می‌کنیم.
	
	برای تصویر ورودی:
	
	\[
	\boxed{
		p_r(r_k)
		=
		\frac{n_k}{MN}
	}
	\tag{19}
	\]
	
	و تابع \lr{Equalization}:
	
	\[
	\boxed{
		s_k
		=
		(L-1)
		\sum_{j=0}^{k}
		p_r(r_j)
	}
	\tag{20}
	\]
	
	است.
	
	اکنون \lr{Histogram} مطلوب را به صورت زیر تعریف می‌کنیم:
	
	\[
	p_z(z_i),
	\qquad
	i=0,1,\ldots,L-1.
	\]
	
	برای آن تابع:
	
	\[
	\boxed{
		G(z_q)
		=
		(L-1)
		\sum_{i=0}^{q}
		p_z(z_i)
	}
	\tag{21}
	\]
	
	محاسبه می‌شود.
	
	
	% =========================================================
	\section{شرط تطبیق $G(z_q)$ با $s_k$}
	
	برای هر مقدار $s_k$ باید مقداری از $z_q$ را پیدا کنیم که:
	
	\[
	\boxed{
		G(z_q)\approx s_k
	}
	\tag{22}
	\]
	
	یعنی:
	
	\[
	\boxed{
		z_q=G^{-1}(s_k)
	}
	\tag{23}
	\]
	
	اما در حالت گسسته معمولاً نیازی به محاسبه تحلیلی $G^{-1}$ نیست.
	
	تمام مقادیر $G(z_q)$ را محاسبه می‌کنیم و نزدیک‌ترین مقدار به $s_k$
	را انتخاب می‌کنیم.
	
	
	% =========================================================
	\section{\lr{Lookup Table}}
	
	برای:
	
	\[
	q=0,1,\ldots,L-1
	\]
	
	تمام مقادیر زیر را محاسبه می‌کنیم:
	
	\[
	G(z_q)
	=
	(L-1)
	\sum_{i=0}^{q}p_z(z_i).
	\]
	
	سپس آنها را گرد می‌کنیم:
	
	\[
	G(z_q)\rightarrow
	\text{\lr{round}}(G(z_q)).
	\]
	
	نتایج در یک:
	
	\[
	\boxed{\text{\lr{Lookup Table (LUT)}}}
	\]
	
	ذخیره می‌شوند.
	
	برای هر $s_k$، نزدیک‌ترین مقدار جدول پیدا می‌شود.
	
	
	% =========================================================
	\section{قاعده انتخاب نزدیک‌ترین مقدار}
	
	برای هر $s_k$، مقدار $z_q$ را طوری انتخاب می‌کنیم که:
	
	\[
	\boxed{
		|G(z_q)-s_k|
	}
	\]
	
	کمینه شود.
	
	یعنی:
	
	\[
	\boxed{
		q^*
		=
		\arg\min_q
		|G(z_q)-s_k|
	}
	\tag{24}
	\]
	
	سپس:
	
	\[
	\boxed{
		s_k\rightarrow z_{q^*}
	}
	\tag{25}
	\]
	
	اگر چند مقدار $z_q$ فاصله برابر داشته باشند، طبق قرارداد کتاب،
	\textbf{کوچک‌ترین مقدار} انتخاب می‌شود.
	
	
	% =========================================================
	\section{الگوریتم کامل \lr{Histogram Specification} گسسته}
	
	\subsection*{مرحله 1: \lr{Equalization} ورودی}
	
	ابتدا \lr{Histogram} تصویر ورودی را محاسبه کنید:
	
	\[
	p_r(r_k)=\frac{n_k}{MN}.
	\]
	
	سپس:
	
	\[
	s_k
	=
	(L-1)
	\sum_{j=0}^{k}p_r(r_j).
	\]
	
	نتایج را به نزدیک‌ترین عدد صحیح گرد کنید.
	
	\subsection*{مرحله 2: \lr{Histogram} مطلوب}
	
	\lr{Histogram} مطلوب را داشته باشید:
	
	\[
	p_z(z_i).
	\]
	
	سپس برای همه $q$:
	
	\[
	G(z_q)
	=
	(L-1)
	\sum_{i=0}^{q}p_z(z_i).
	\]
	
	مقادیر را گرد کرده و در \lr{LUT} قرار دهید.
	
	\subsection*{مرحله 3: یافتن نگاشت}
	
	برای هر $s_k$ مقدار $z_q$ را پیدا کنید که:
	
	\[
	|G(z_q)-s_k|
	\]
	
	کمترین مقدار را داشته باشد.
	
	سپس:
	
	\[
	s_k\rightarrow z_q.
	\]
	
	\subsection*{مرحله 4: ساخت تصویر نهایی}
	
	هر پیکسل تصویر \lr{Equalized} که مقدار آن $s_k$ است، به مقدار متناظر
	$z_q$ تبدیل می‌شود:
	
	\[
	\boxed{
		s_k\rightarrow z_q
	}
	\]
	
	در نهایت تصویر دارای \lr{Histogram} نزدیک به \lr{Histogram} مشخص‌شده خواهد بود.
	
	
	% =========================================================
	\section{مثال \lr{Histogram Specification}}
	
	فرض کنید تصویر ورودی همان تصویر 3 بیتی مثال قبل باشد:
	
	\[
	L=8.
	\]
	
	پس مقادیر \lr{Equalized} به صورت تقریبی:
	
	\[
	\boxed{
		s_0,s_1,\ldots,s_7
		=
		1,3,5,6,7,7,7,7
	}
	\tag{26}
	\]
	
	فرض کنید \lr{Histogram} مطلوب به صورت $p_z(z_q)$ داده شده باشد.
	
	با استفاده از رابطه:
	
	\[
	G(z_q)
	=
	7
	\sum_{i=0}^{q}p_z(z_i)
	\]
	
	مقادیر $G$ را محاسبه می‌کنیم.
	
	پس از گرد کردن، ممکن است جدولی مانند جدول زیر حاصل شود:
	
	\[
	\begin{array}{c|c}
		z_q&G(z_q)\\
		\hline
		0&0\\
		1&0\\
		2&2\\
		3&1\\
		4&5\\
		5&6\\
		6&6\\
		7&7
	\end{array}
	\]
	
	در عمل برای هر $s_k$ نزدیک‌ترین $G(z_q)$ انتخاب می‌شود.
	
	
	% =========================================================
	\section{مثال نگاشت $s$ به $z$}
	
	اگر:
	
	\[
	s_0=1
	\]
	
	باشد، بررسی می‌کنیم که کدام مقدار $G(z_q)$ به 1 نزدیک‌تر است.
	
	اگر:
	
	\[
	G(z_3)=1,
	\]
	
	آنگاه:
	
	\[
	\boxed{
		s_0=1\rightarrow z_3=3
	}
	\]
	
	یعنی تمام پیکسل‌هایی که در تصویر \lr{Equalized} دارای شدت 1 هستند،
	در تصویر نهایی دارای شدت 3 خواهند شد.
	
	اگر مثلاً 790 پیکسل مقدار 1 داشته باشند، در تصویر نهایی:
	
	\[
	n_z(3)=790.
	\]
	
	پس احتمال شدت 3 برابر است با:
	
	\[
	p_z(3)
	=
	\frac{790}{4096}.
	\]
	
	یعنی:
	
	\[
	\boxed{
		p_z(3)\approx0.19.
	}
	\]
	
	
	% =========================================================
	\section{چرا \lr{Histogram} نهایی دقیقاً برابر \lr{Histogram} مطلوب نیست؟}
	
	در حالت پیوسته می‌توان نتیجه دقیق را به دست آورد.
	
	اما در حالت دیجیتال:
	
	\begin{enumerate}
		\item تعداد شدت‌ها محدود است؛
		\item تعداد پیکسل‌ها محدود است؛
		\item $s_k$ باید گرد شود؛
		\item $G(z_q)$ نیز باید گرد شود؛
		\item چند مقدار $s_k$ ممکن است به یک $z_q$ نگاشت شوند.
	\end{enumerate}
	
	بنابراین معمولاً:
	
	\[
	\boxed{
		\text{\lr{Histogram}}_{\text{\lr{output}}}
		\ne
		\text{\lr{Histogram}}_{\text{\lr{specified}}}
	}
	\]
	
	اما شکل کلی توزیع می‌تواند به \lr{Histogram} موردنظر نزدیک شود.
	
	
	% =========================================================
	\section{شرط \lr{Strict Monotonicity} در \lr{Histogram Specification}}
	
	برای اینکه $G^{-1}$ به صورت یکتا وجود داشته باشد، باید $G$ کاملاً
	صعودی باشد.
	
	از رابطه:
	
	\[
	G(z_q)
	=
	(L-1)
	\sum_{i=0}^{q}p_z(z_i)
	\]
	
	اگر یکی از احتمال‌ها صفر باشد:
	
	\[
	p_z(z_i)=0,
	\]
	
	ممکن است داشته باشیم:
	
	\[
	G(z_i)=G(z_{i+1}).
	\]
	
	در این صورت $G$ دیگر \lr{Strictly Monotonic} نیست.
	
	بنابراین برای داشتن تابع معکوس کاملاً یکتا، لازم است:
	
	\[
	\boxed{
		p_z(z_i)>0
		\qquad
		\forall i.
	}
	\tag{27}
	\]
	
	اما در عمل اگر برخی مقادیر صفر باشند، از روش نزدیک‌ترین مقدار در
	\lr{Lookup Table} استفاده می‌کنیم.
	
	
	% =========================================================
	\section{مقایسه \lr{Histogram Equalization} و \lr{Histogram Specification}}
	
	\begin{center}
		\begin{tabular}{>{\raggedleft\arraybackslash}p{4cm}
				>{\raggedleft\arraybackslash}p{5cm}
				>{\raggedleft\arraybackslash}p{5cm}}
			\toprule
			ویژگی & \lr{Histogram Equalization} & \lr{Histogram Specification}\\
			\midrule
			هدف &
			افزایش کنتراست &
			رسیدن به \lr{Histogram} مطلوب\\
			\hline
			\lr{Histogram} هدف &
			تقریباً \lr{Uniform} &
			از قبل مشخص می‌شود\\
			\hline
			نیاز به \lr{Histogram} مطلوب &
			خیر &
			بله\\
			\hline
			پیچیدگی &
			ساده &
			بیشتر\\
			\hline
			کاربرد &
			بهبود خودکار کنتراست &
			کنترل شکل \lr{Histogram} خروجی\\
			\bottomrule
		\end{tabular}
	\end{center}
	
	
	% =========================================================
	\section{خلاصه ریاضی کل مبحث}
	
	\subsection{\lr{Histogram} تصویر}
	
	\[
	\boxed{
		p_r(r_k)=\frac{n_k}{MN}
	}
	\]
	
	\subsection{\lr{Histogram Equalization} پیوسته}
	
	\[
	\boxed{
		s=T(r)
		=
		(L-1)
		\int_0^r p_r(w)\,dw
	}
	\]
	
	\subsection{رابطه تبدیل \lr{PDF}}
	
	\[
	\boxed{
		p_s(s)
		=
		p_r(r)
		\left|
		\frac{dr}{ds}
		\right|
	}
	\]
	
	\subsection{نتیجه \lr{Equalization}}
	
	\[
	\boxed{
		p_s(s)=\frac{1}{L-1}
	}
	\]
	
	\subsection{\lr{Histogram Equalization} گسسته}
	
	\[
	\boxed{
		s_k
		=
		(L-1)
		\sum_{j=0}^{k}p_r(r_j)
	}
	\]
	
	و:
	
	\[
	\boxed{
		s_k
		=
		(L-1)
		\sum_{j=0}^{k}
		\frac{n_j}{MN}
	}
	\]
	
	\subsection{\lr{Histogram} مطلوب}
	
	\[
	\boxed{
		G(z)
		=
		(L-1)
		\int_0^z p_z(v)\,dv
	}
	\]
	
	\subsection{رابطه اساسی \lr{Histogram Specification}}
	
	\[
	\boxed{
		G(z)=T(r)
	}
	\]
	
	بنابراین:
	
	\[
	\boxed{
		z=G^{-1}(s)
	}
	\]
	
	و مستقیماً:
	
	\[
	\boxed{
		z=G^{-1}[T(r)]
	}
	\]
	
	\subsection{فرم گسسته تابع $G$}
	
	\[
	\boxed{
		G(z_q)
		=
		(L-1)
		\sum_{i=0}^{q}p_z(z_i)
	}
	\]
	
	\subsection{انتخاب مقدار متناظر}
	
	\[
	\boxed{
		q^*
		=
		\arg\min_q
		|G(z_q)-s_k|
	}
	\]
	
	و:
	
	\[
	\boxed{
		s_k\rightarrow z_{q^*}.
	}
	\]
	
	
	% =========================================================
	\section{درک شهودی کل فرآیند}
	
	می‌توان کل مبحث را در دو زنجیره بسیار ساده خلاصه کرد.
	
	\subsection*{\lr{Histogram Equalization}}
	
	\[
	\boxed{
		r
		\longrightarrow
		\text{\lr{CDF}}
		\longrightarrow
		s
	}
	\]
	
	یعنی:
	
	\[
	\boxed{
		s=(L-1)\text{\lr{CDF}}(r)
	}
	\]
	
	هدف:
	
	\[
	\boxed{
		\text{گسترش کنتراست}
	}
	\]
	
	\subsection*{\lr{Histogram Specification}}
	
	\[
	\boxed{
		r
		\longrightarrow
		s
		\longrightarrow
		z
	}
	\]
	
	که:
	
	\[
	s=T(r)
	\]
	
	و:
	
	\[
	z=G^{-1}(s).
	\]
	
	در نتیجه:
	
	\[
	\boxed{
		z=G^{-1}[T(r)]
	}
	\]
	
	هدف:
	
	\[
	\boxed{
		\text{رسیدن به یک \lr{Histogram} مشخص}
	}
	\]
	
	
	% =========================================================
	

% =========================================================
\section{بسته کد همراه و اجرای یک‌مرحله‌ای}

همراه این جزوه پوشه \lr{code} ارائه شده است. کدها برای \lr{Python 3.12} نوشته شده‌اند
و نسخه همراه جزوه به‌صورت آفلاین و مستقل از مخازن اینترنتی طراحی شده است. در اجرای
عادی هیچ فرمان \lr{pip install} اجرا نمی‌شود، به \lr{PyPI} اتصال برقرار نمی‌شود و
کتابخانه‌های ثالثی مانند \lr{NumPy}، \lr{Matplotlib} یا \lr{OpenCV} لازم نیستند.
اگر تصویر ورودی داده نشود، برنامه یک تصویر آزمایشی مصنوعی می‌سازد؛ بنابراین تمام
آزمایش‌های آموزشی تنها با کتابخانه استاندارد \lr{Python 3.12} اجرا می‌شوند.

موضوع‌های کدنویسی شامل محاسبه هیستوگرام نرمال‌شده، متعادل‌سازی سراسری با ساخت مستقیم
\lr{LUT} از \lr{CDF}، تطبیق هیستوگرام، \lr{CLAHE}، روش \lr{BBHE}، آنتروپی،
واگرایی \lr{KL} نسبت به توزیع یکنواخت، هیستوگرام مشترک و اطلاعات متقابل است. خروجی‌ها
در پوشه \lr{outputs} ذخیره می‌شوند و فایل‌های \lr{CSV} و \lr{metrics.txt} مقادیر
عددی قبل و بعد از تبدیل را گزارش می‌کنند. تصاویر خروجی \lr{PNG} نیز مستقیماً با
امکانات استاندارد زبان تولید می‌شوند.

برای اجرای مستقیم تنها کافی است در ویندوز روی

\[
\boxed{\text{\lr{RUN\_ALL.bat}}}
\]

دوبار کلیک شود. این فایل فقط وجود نسخه ۶۴ بیتی \lr{Python 3.12} را بررسی می‌کند و سپس
همه مثال‌ها را اجرا می‌کند؛ محیط مجازی و فایل نیازمندی برای اجرای اصلی لازم نیست.
پیاده‌سازی متعادل‌سازی سراسری، تطبیق هیستوگرام و \lr{CLAHE} در این بسته عمداً از روی
روابط و الگوریتم‌های پایه نوشته شده است تا ارتباط میان معادلات جزوه و کد قابل مشاهده
باشد. در پیاده‌سازی \lr{CLAHE}، هیستوگرام هر ناحیه برش داده می‌شود، جرم اضافی دوباره
میان سطوح خاکستری توزیع می‌شود و نگاشت‌های محلی با درون‌یابی دوخطی ترکیب می‌شوند.

در حالت کاملاً مستقل از کتابخانه‌های جانبی، ورودی‌های \lr{PGM} و \lr{BMP} پشتیبانی
می‌شوند. اگر کتابخانه \lr{Pillow} از قبل روی سیستم موجود باشد، خواندن \lr{PNG} و
\lr{JPEG} نیز به‌عنوان یک قابلیت اختیاری فعال می‌شود؛ با این حال هیچ‌یک از مثال‌های
اصلی جزوه به آن وابسته نیستند.


\section{نکته بسیار مهم امتحانی}
	
	سه رابطه زیر را باید به خوبی به خاطر سپرد:
	
	\[
	\boxed{
		p_r(r_k)=\frac{n_k}{MN}
	}
	\]
	
	\[
	\boxed{
		s_k=(L-1)\sum_{j=0}^{k}p_r(r_j)
	}
	\]
	
	و برای \lr{Histogram Specification}:
	
	\[
	\boxed{
		G(z_q)
		=
		(L-1)
		\sum_{i=0}^{q}p_z(z_i)
	}
	\]
	
	سپس:
	
	\[
	\boxed{
		s_k
		\rightarrow
		z_q
		\quad
		\text{با انتخاب نزدیک‌ترین }
		G(z_q)
		\text{ به }
		s_k.
	}
	\]
	
	به بیان مفهومی:
	
	\[
	\boxed{
		\text{\lr{Histogram Equalization}}
		=
		\text{\lr{CDF} ورودی}
	}
	\]
	
	و:
	
	\[
	\boxed{
		\text{\lr{Histogram Specification}}
		=
		\text{\lr{Equalization}}
		+
		\text{\lr{Inverse CDF} مطلوب}.
	}
	\]
	


% =========================================================
\section{منابع}
\begin{thebibliography}{99}

\bibitem{pizer1987}
\lr{S. M. Pizer, E. P. Amburn, J. D. Austin, R. Cromartie, A. Geselowitz, T. Greer, B. ter Haar Romeny, J. B. Zimmerman, and K. Zuiderveld, ``Adaptive histogram equalization and its variations,'' Computer Vision, Graphics, and Image Processing, vol. 39, no. 3, pp. 355--368, 1987.}
\href{https://doi.org/10.1016/S0734-189X(87)80186-X}{\lr{DOI: 10.1016/S0734-189X(87)80186-X}}

\bibitem{kim1997}
\lr{Y.-T. Kim, ``Contrast enhancement using brightness preserving bi-histogram equalization,'' IEEE Transactions on Consumer Electronics, vol. 43, no. 1, pp. 1--8, 1997.}
\href{https://doi.org/10.1109/30.580378}{\lr{DOI: 10.1109/30.580378}}

\bibitem{stark2000}
\lr{J. A. Stark, ``Adaptive image contrast enhancement using generalizations of histogram equalization,'' IEEE Transactions on Image Processing, vol. 9, no. 5, pp. 889--896, 2000.}
\href{https://doi.org/10.1109/83.841534}{\lr{DOI: 10.1109/83.841534}}

\bibitem{arici2009}
\lr{T. Arici, S. Dikbas, and Y. Altunbasak, ``A Histogram Modification Framework and Its Application for Image Contrast Enhancement,'' IEEE Transactions on Image Processing, vol. 18, no. 9, pp. 1921--1935, 2009.}
\href{https://doi.org/10.1109/TIP.2009.2021548}{\lr{DOI: 10.1109/TIP.2009.2021548}}

\bibitem{shannon1948}
\lr{C. E. Shannon, ``A Mathematical Theory of Communication,'' Bell System Technical Journal, vol. 27, no. 3, pp. 379--423, 1948.}
\href{https://doi.org/10.1002/j.1538-7305.1948.tb01338.x}{\lr{DOI: 10.1002/j.1538-7305.1948.tb01338.x}}

\bibitem{kullback1951}
\lr{S. Kullback and R. A. Leibler, ``On Information and Sufficiency,'' The Annals of Mathematical Statistics, vol. 22, no. 1, pp. 79--86, 1951.}
\href{https://doi.org/10.1214/aoms/1177729694}{\lr{DOI: 10.1214/aoms/1177729694}}

\bibitem{celik2014}
\lr{T. Celik, ``Spatial Entropy-Based Global and Local Image Contrast Enhancement,'' IEEE Transactions on Image Processing, vol. 23, no. 12, pp. 5298--5308, 2014.}
\href{https://doi.org/10.1109/TIP.2014.2364537}{\lr{DOI: 10.1109/TIP.2014.2364537}}

\bibitem{zhuang2018}
\lr{L. Zhuang and Y. Guan, ``Adaptive Image Enhancement Using Entropy-Based Subhistogram Equalization,'' Computational Intelligence and Neuroscience, vol. 2018, Article ID 3837275, 2018.}
\href{https://doi.org/10.1155/2018/3837275}{\lr{DOI: 10.1155/2018/3837275}}

\bibitem{viola1997}
\lr{P. Viola and W. M. Wells III, ``Alignment by Maximization of Mutual Information,'' International Journal of Computer Vision, vol. 24, no. 2, pp. 137--154, 1997.}
\href{https://doi.org/10.1023/A:1007958904918}{\lr{DOI: 10.1023/A:1007958904918}}

\end{thebibliography}

\end{document}
